% ===========================================================================
% P1.10 DRAFT - new/rewritten sections for main_en.tex (Khung A framing).
% Written 2026-10-07. NOT yet merged into main_en.tex: pending author
% decisions C2 (title), C4 (journal), C5 (force_periodic). Blocks marked
% [C5] change wording if periodic-edge enforcement is dropped.
% Every number is traceable to outputs/phase5/plan_v3/ (source in comments).
% Uses the macros of main_en.tex (\nuot, \nuto, \RR).
% ===========================================================================

% ---------------------------------------------------------------------------
% REPLACES the abstract (keeps the "optimize what is verified" core; removes
% the multi-objective/OOD claims that P1.6a-b, P1.7 do not support).
% ---------------------------------------------------------------------------
\begin{abstract}
Gradient-based design of metamaterial unit cells, whether by topology
optimization or by searching the latent space of a generative model, is
judged by finite-element (FE) analysis of a binarized design, yet it usually
optimizes a continuous density field. We show that this mismatch is
systematically exploited. Refining the latent code of a conditional
variational autoencoder (cVAE) against exact homogenization sensitivities
drives the objective to $\sim\!10^{-10}$ while the verified error stays six
orders of magnitude larger, and density-based topology optimization with
Heaviside continuation to $\beta=64$ shows the same gap (median $2\times10^{-7}$
before versus $3.6\times10^{-2}$ after thresholding). Placing the verification
operators inside the refinement objective raises the FE-verified $\RR$ of
the Poisson ratio $\nuot$ of two-dimensional auxetic cells from 0.874 to 0.998
for single samples (0.831 to 0.990 on a second, disjoint target set) and
from 0.983--0.990 to 0.995--0.998 after best-of-30 selection on three sets.
Against topology optimization from scratch with multi-start, the generative
pipeline reaches the same accuracy with about eight times fewer FE solves
and, on a pre-registered test set, a 27\% lower joint error for targets of
moderate anisotropy. It fails for strongly anisotropic targets that are
sparse in the training data, and topology optimization remains superior in
manufacturability and for targets outside the training range.
\end{abstract}

% ---------------------------------------------------------------------------
% NEW Results subsection (after "Verified accuracy requires optimizing the
% verified design"; before multi-objective).
% Sources: p1_6a_simp_baseline_in100.json, p1_7_simp_in100b_full.json,
% p1_8_simp_in100c_full.json, p1_6c/p1_7_in100b/p1_8_in100c_bo30_images.json,
% p1_8_anisotropy_strata.json, p1_7_hybrid_in100.json, p1_6x_*.json.
% ---------------------------------------------------------------------------
\subsection{Comparison with topology optimization from scratch}

The relevant competitor of a generative pipeline is not a lookup table but
topology optimization run directly for the target. We solved the classical
inverse homogenization problem $\min\tfrac12\lVert\boldsymbol\nu-\boldsymbol\nu^{*}\rVert^2$
for every target with the method of moving asymptotes, a volume bound and
stiffness floors, Heaviside continuation to $\beta=64$, and four starting
topologies taken from the training data (Methods). Each run was verified
exactly like a generated design, and the best of the four was selected by
its verified error. This baseline uses on average 893 FE solves per target,
against 107 for best-of-30 sampling followed by guarded binarization-aware
refinement.

Topology optimization exhibits the same objective--verification gap as
latent refinement. At the end of continuation the median error on the
projected density field was $2\times10^{-7}$, whereas the median verified
error after thresholding was $3.6\times10^{-2}$ per run, and for 57\% of the
300 targets the starting topology with the lowest projected error was not
the one with the lowest verified error. Selection by verified error is therefore
necessary for both methods, which supports the central claim of this paper
beyond generative models.

Table~\ref{tab:simp} compares the two methods on three disjoint sets of 100
held-out targets. On all three, the generative pipeline matches converged
topology optimization in $\nuot$ ($\RR=0.994$--$0.998$ versus
$0.996$--$0.998$) with roughly eightfold fewer FE solves. Over all targets
its joint error $|\Delta\nuot|+|\Delta\nuto|$ is not consistently lower,
because of a specific failure mode. Every disconnected cell produced by the
cVAE (one target in the first set, two in the second) belongs to a target
with an anisotropy ratio $r=\max(\nuto^{*}/\nuot^{*},\,\nuot^{*}/\nuto^{*})\ge10$,
which corresponds to a stiffness ratio $E_2/E_1\ge10$ for an orthotropic
cell. For these targets all 30 sampled candidates were disconnected, so the
failure cannot be fixed by selection; only 2.8\% of the training samples
have $r>10$. Having observed this on the first two sets, we pre-registered a
stratified test on a third set before running it: for targets with
$r<10$, the generative pipeline should have a lower joint error and a lower
$\nuot$ error than topology optimization, with paired 95\% CIs excluding zero.
The joint-error criterion was met (27\% lower, CI 8--42\%); the $\nuot$
criterion was not (19\% lower, CI $-10$ to 40\%). This test used the
pipeline with periodic-edge enforcement. After that step was removed for
conceptual reasons (Methods), the same comparison favours the generative
pipeline in both criteria on all three sets (Table~\ref{tab:simp}); because
this version was not pre-registered, we report it as a sensitivity analysis
and still claim parity in $\nuot$ and an advantage in joint error for
moderate anisotropy, not superiority in $\nuot$. The differences in $\nuot$ between the two methods
($\sim\!10^{-3}$) are in any case an order of magnitude smaller than the
discretization error of the $50\times50$ FE model (next paragraph).

\begin{table}[t]
\centering
\caption{Generative pipeline (best-of-30 + guarded binarization-aware
refinement) versus topology optimization from scratch (best of four
starts, converged) on three disjoint sets of 100 held-out targets. Both
methods select candidates by verified joint error (not by the composite
score of Table~\ref{tab:multi}). FE:
solves per target. Manufacturable: single connected solid phase with
minimum feature size of two pixels. Stratified rows: relative reduction of
the error of the generative pipeline versus topology optimization for targets
with anisotropy ratio $r<10$ (paired bootstrap 95\% CI; positive favours the
generative pipeline). Final pipeline without periodic-edge enforcement; the
pre-registered test on set C was run with it (see text; both versions in the
Supplementary Information).}
\label{tab:simp}
\small
\setlength{\tabcolsep}{3.5pt}
\begin{tabular}{llccccc}
\toprule
Set & Method & $\RR(\nuot)$ & $\RR(\nuto)$ & Joint MAE & FE & Manuf.\\
\midrule
\multirow{2}{*}{A} & generative & 0.998 & 0.846 & 0.0139 & 107 & 0.32\\
 & topology opt. & 0.997 & 0.997 & 0.0128 & 897 & 0.77\\
\multirow{2}{*}{B} & generative & 0.995 & 0.745 & 0.0229 & 107 & 0.35\\
 & topology opt. & 0.996 & 0.997 & 0.0139 & 889 & 0.74\\
\multirow{2}{*}{C} & generative & 0.998 & 0.998 & 0.0082 & 107 & 0.32\\
 & topology opt. & 0.998 & 0.996 & 0.0126 & 893 & 0.83\\
\midrule
\multicolumn{7}{l}{$r<10$: joint error \;A $+41\%$ [26, 54] \;B $+40\%$ [23, 53] \;C $+35\%$ [19, 48]}\\
\multicolumn{7}{l}{$r<10$: $\nuot$ error \;A $+39\%$ [19, 54] \;B $+39\%$ [19, 54] \;C $+37\%$ [14, 54]}\\
\bottomrule
\end{tabular}
\end{table}

\paragraph{Mesh sensitivity.}
All accuracies above are relative to the $50\times50$ FE model used for
data generation, refinement and verification. Re-solving the same binary
geometries with each pixel subdivided into $4\times4$ elements changes
$\nuot$ by a median of 0.015 (3.8\%) and at most 0.042 on 30 test designs,
towards more negative values in 27 of them. On the final designs of the
first set the ranking of the two methods is unchanged on the $200\times200$
mesh ($\RR(\nuot)$ 0.992 versus 0.985).

\paragraph{A hybrid does not inherit the strengths of both.}
Because topology optimization was more robust and more manufacturable while
the generative pipeline was cheaper, we pre-registered a hybrid in which the
generated design initializes topology optimization (continuation
$\beta=8\to64$, at most 61 additional FE solves). It met the cost and
non-degeneracy criteria (166 FE solves per target; 1\% degenerate) but
failed the joint-error non-inferiority criterion against converged
topology optimization (CI lower bound $-96\%$, driven by the one
disconnected target) and the manufacturability criterion (0.33 against a
threshold of 0.70). Starting from a sharp projection preserves the topology,
including the thin features, of the generated design.

% ---------------------------------------------------------------------------
% NEW Methods subsection.
% ---------------------------------------------------------------------------
\subsection{Topology optimization baseline}
For each target we minimized
$\tfrac12\lVert\boldsymbol\nu(\hat{\mathbf x})-\boldsymbol\nu^{*}\rVert_2^2$
over element densities $\mathbf x$, with $\hat{\mathbf x}=H_\beta(\tilde{\mathbf x})$
and $\tilde{\mathbf x}$ the density-filtered field (radius 1.5 elements),
subject to $\overline{\hat{\mathbf x}}\le0.55$ and $Q_{11},Q_{22}\ge0.1\cdot0.55\,E_0$,
using the method of moving asymptotes \citep{svanberg1987mma} with penalization
$p=3$. Sensitivities follow from the same analytic expressions used for
refinement and were checked against finite differences. Continuation runs
$\beta\in\{1,2,4,8,16,32,64\}$ with at most 43 evaluations per level and a
restart of the optimizer at each level. Starting topologies are the four
most frequent seed families of the training data, rescaled to the volume
bound. The final $\hat{\mathbf x}$ is thresholded at 0.5 and verified with an
independent FE solve. Parameters are the medians of the training data.
The hybrid uses the generated binary design as the starting point and
$\beta\in\{8,16,32,64\}$ with at most 15 evaluations per level.

% ---------------------------------------------------------------------------
% ADDITION to "Statistics and pre-registration".
% ---------------------------------------------------------------------------
Two further experiments were specified in writing before they were run: the
hybrid (four criteria: joint-error non-inferiority within 10\%, at least 70\%
manufacturable, at most 299 FE solves per target, at most 1\% degenerate
designs) and the stratified comparison on the third target set (both CIs
excluding zero for $r<10$). The first two target sets were analysed before
the stratification was conceived, and their stratified results are reported
as post hoc.

% ---------------------------------------------------------------------------
% REPLACES the Limitations paragraph.
% ---------------------------------------------------------------------------
\paragraph{Limitations.}
(i) Accuracy is measured against a $50\times50$ linear-elastic,
small-strain FE model with a base-material Poisson ratio of 0.3; the
discretization error of this model (median 0.015 in $\nuot$) exceeds the
reported mean errors. Large strain and three-dimensional cells are not
addressed, and no design was tested experimentally.
(ii) The generative pipeline fails for strongly anisotropic targets
($r\ge10$), which are sparse in the training data; topology optimization does
not.
(iii) Topology optimization from scratch produces manufacturable cells more than
twice as often (0.74--0.83 versus 0.32--0.35). Refinement does not
change this systematically: across five runs the manufacturable fraction
moved by $-6$ to $+9$ percentage points (61 designs became manufacturable,
44 ceased to be; sign test $p=0.12$).
(iv) For non-auxetic targets outside the training range, topology
optimization is markedly more accurate (MAE in $\nuot$ 0.005--0.040 versus
0.06--0.09); the advantage of the generative model in that regime holds only
against retrieval from the training library.
(v) The cost advantage is amortized: generating the training data required
about $2.7$--$3.0\times10^{6}$ FE solves, which pays off against converged
topology optimization only after about 3\,500 targets.
(vi) The aesthetic score is a heuristic that has not been validated against
human judgement.

% ---------------------------------------------------------------------------
% REPLACES Discussion paragraphs "Accuracy", "Multiple objectives" and
% "Out-of-distribution generalization" (Limitations: see block above).
% [C5] marks wording that changes if periodic-edge enforcement is dropped.
% ---------------------------------------------------------------------------
\paragraph{Optimizing what is verified.}
The central finding is that any gradient-based search over density fields,
whether in the latent space of a generative model or over element densities
in topology optimization, will exploit the difference between the field it
differentiates and the binarized design that is finally analysed. The
optimizer reaches objectives of $10^{-10}$ on fields that the verifier sees
very differently (Fig.~\ref{fig:gap}), and the same happens in topology
optimization with Heaviside continuation to $\beta=64$. The gap is invisible
in the loss curve. Two remedies follow, and both are cheap: include the
verification operators in the objective, and select among candidates by
verified rather than optimized error. Neither ingredient is new in
density-based topology optimization \citep{guest2004achieving,wang2011projection};
what this work shows is that omitting them changes the conclusions of a
comparison between methods, and that the generative and the classical route
need them equally.

\paragraph{What a generative model adds.}
Measured against converged topology optimization rather than against a
lookup table, the advantage of the generative pipeline is narrower than is
commonly implied. It reaches the same verified accuracy with about eight
times fewer FE solves per target, and a lower joint error for targets of
moderate anisotropy, but the saving is amortized only after thousands of
targets because the training data themselves were produced by topology
optimization. The pipeline is therefore attractive when many targets must be
served, for example in interactive design or in the inner loop of a
multiscale optimization, and not when a single cell is needed. Its failures
are predictable: strongly anisotropic targets that are sparse in the training
data. Reporting such failures by stratum, rather than averaging them into a
single score, made them visible in our study and should become standard.

\paragraph{Manufacturability and out-of-distribution targets.}
Topology optimization from scratch produced connected cells with adequate
feature size more than twice as often as the generative pipeline, and
latent refinement did not change that fraction systematically. Initializing topology
optimization with the generated design did not recover this property,
because a sharp initial projection preserves the generated topology. For
non-auxetic targets outside the training range, the cVAE reaches the correct
sign where retrieval from the library cannot, but topology optimization
reaches the target magnitude far more accurately. Along magnitude, all
generative variants stop at the edge of the training data. OOD claims for
generative models should therefore name the direction in property space,
check that the target is physically admissible, and compare against
topology optimization rather than retrieval. [C5: if periodic-edge
enforcement is dropped, add one sentence that element-based periodic
boundary conditions make every pixel cell tile by translation, so no edge
matching is required.]
